% 連続体力学 連続体力学-1 問題3（003-kijutsu） 解いた日 2026-10-17
%% kotae
\toi{1}{○}{記述}
質点…大きさをもたない。

剛体…原子・分子の相対位置が変わらない。

連続体…原子・分子の相対位置が変わる。

歪み…原子・分子の相対位置\nuke{の変化}
\naze{歪みは相対位置そのものではなく，相対位置の「変化」である。}
\toi{2}{○}{記述}
固定：Euler

動く：Lagrange

移流項：$\vect{u}\cdot\vect{\nabla}$\nuke{$(\vect{u}\cdot\vect{\nabla})T$}
\naze{問いは差を表す項の式を求めている。演算子だけでなく，$\mathrm{D}T/\mathrm{D}t-\partial T/\partial t=(\vect{u}\cdot\vect{\nabla})T=\sum_{i}u_{i}\,\partial T/\partial x_{i}$ と $T$ に作用させた形で書く。}
\toi{3}{○}{記述}
ある。

水温の分布に沿って流れる粒子の温度が変わるから。
\naze{理由が結論の言い換えになっている。$\partial T/\partial t=0$ でも $\mathrm{D}T/\mathrm{D}t=(\vect{u}\cdot\vect{\nabla})T$ が残り，流れの向きに温度の勾配があれば0でない，と式で根拠を書く。}
%% kizuki
\item 歪みは原子・分子の相対位置の「変化」である。剛体（相対位置が固定）と連続体（相対位置が変化）の違いがそのまま歪みの定義になっている。
\item 2つの温度計の差は $\mathrm{D}T/\mathrm{D}t-\partial T/\partial t=(\vect{u}\cdot\vect{\nabla})T$ で，(3) はこの式で $\partial T/\partial t=0$ とおけば移流項だけが残る，と答えられる。「式を答えよ」「理由とともに」では，名前や言い換えで済ませず式を1行書く。
